\section{Positioning: fdtdKit and other electromagnetic solvers}\label{sec:positioning}
fdtdKit makes one open-source solver faster on one class of machine; it does not widen what openEMS can model. This
section places it against the tools an RF engineer would otherwise use, so that the speed-ups of
Section~\ref{sec:results} are not read as a claim they do not make. No commercial solver was run for this report;
the comparisons with them below are qualitative or, where a number appears, a roofline bound~\cite{roofline} computed
from published hardware figures, not a measurement.

\paragraph{Commercial solvers.} Industrial RF design runs mostly on commercial packages. Ansys HFSS solves
Maxwell's equations in the frequency domain with the finite-element method on an automatically and adaptively
refined tetrahedral mesh~\cite{hfssmesh}; this suits high-$Q$ resonators, curved geometry and fine features that a
Cartesian staircase mesh resolves poorly. CST Studio Suite offers, among several solvers, a time-domain
finite-integration solver closely related to FDTD, with GPU acceleration~\cite{cstgpu}; Remcom XFdtd is an FDTD
solver with GPU acceleration as well~\cite{xfdtd}. Beyond the solvers, these packages bring CAD import, automatic
meshing with convergence control, parameter optimisation, a graphical front end, vendor validation and support.
fdtdKit touches none of this: model setup, meshing and post-processing are openEMS's and the author's scripts.

\paragraph{Raw speed.} Section~\ref{sec:ceiling} shows that the CPU engine runs close to what the cores can draw from
memory. A commercial FDTD/FIT solver on a data-centre GPU is therefore bounded far higher: an NVIDIA H100 SXM has
3.35\,TB/s of memory bandwidth~\cite{h100}, which at the 74 bytes per cell of an unfused update is a roofline bound of
about 45\,000\,\MC, some 20 times the M5's bound for the same traffic (Section~\ref{sec:ceiling}). Real codes reach a
fraction of their roofline, and we have no measurement of any commercial code, but on large models fdtdKit should not
be expected to match such a system. Its comparison class is a laptop or desktop without a licence.

\paragraph{Open-source solvers.} Among free tools, MEEP~\cite{meep} is the reference FDTD code for photonics and
gprMax~\cite{gprmax} for ground-penetrating radar; openEMS is the usual choice for PCB-scale RF and antenna work,
largely for its Matlab/Octave and Python interfaces. fdtdKit keeps openEMS's interfaces unchanged and its results bit
for bit identical to stock for the same number of time steps (Section~\ref{sec:method}), so existing scripts gain the
speed-up without modification (for SAR the two Python modules must be rebuilt against the new libraries,
\code{BUILD.md} section 6). That holds for the CPU engines, which stay the default. The opt-in Metal engine
(\code{--engine=metal-diamond}, Section~\ref{sec:gpu}) is bit-identical to the multithreaded CPU engine run with
\code{OPENEMS\_FTZ=1}, not to stock, because Apple's GPU flushes subnormal floats to zero; against stock its port
probes differ by $1.7\times10^{-7}$ to
$9.4\times10^{-7}$ relative (Section~\ref{sec:policy}). A script that ends on the energy end criterion can stop at a
different step than stock v0.37.0, because the fork's base checks that criterion on a different schedule (upstream
commit eb1c5ee, Section~\ref{sec:policy}); for the same number of steps the records are identical.

\paragraph{Prior work on FDTD and openEMS performance.} That low-order FDTD is bound by memory bandwidth on
multicore CPUs, and that the remedy is to reuse data in cache, is well established: Cessenat reports a more than
twofold reduction from cache-aware restructuring of the Sophie code~\cite{sophie}, and temporal (time-space) tiling
of FDTD and of stencils in general has been studied extensively~\cite{fukaya2018,malas2015}. For openEMS itself, an
openEMS contributor (GitHub user \code{biergaizi}) has documented the engine's synchronisation and memory-traffic
bottlenecks and developed an experimental tiling engine, reported bit-identical to the multithreaded engine, with
speed-ups of about 2--6$\times$ reported in the project's discussion forum~\cite{biergaizi154,biergaizi29}; at the
time of writing that work is not merged. fdtdKit's CPU engine reaches a comparable range by a different route (fewer
barriers, a leaner inner loop, a fused single-pass sweep, the allocation stagger and the dispersive-material hooks)
without temporal tiling. Two time steps per sweep gained 10--14\,\% in our first standalone kernel, which ran the stock
arithmetic (invariants hoisted, the 12 coefficient tables) as a fused sweep on one independent sub-grid per thread
(project log; Section~\ref{sec:rejected}); with the current engine's loops in the same set-up, trapezoidal blocking of
two to four steps per sweep (after MrHighVoltage's fork, next paragraph) gained at most 6\,\% as an upper bound, measured
without halo lines or extension hooks (\code{results/cpu\_items\_cg29\_2026-11-05.txt}, section 2), so the CPU engine's
default sweep stays unblocked. Upstream master's new \code{cpu} engine
blocks with one tile for all threads (Section~\ref{sec:threebuilds} has what that is worth on this machine), and
since 2026-11-11 the fork has that form as an opt-in tiled sweep, which at 4 threads is ahead of the fused sweep at
its best on two large grids and behind on four other models (Section~\ref{sec:tiled}). Wavefront or diamond blocking
beyond that sweep's
trapezoids along x was not tried on the CPU.
The Metal engine of
Section~\ref{sec:gpu} does block over time steps, in diamond tiles on the schedule of another contributor's Metal
engine (next paragraph). Both fdtdKit's CPU factors and the 2--6$\times$ above are measured against stock openEMS, whose
sweep is unfused, and the two approaches are not exclusive. The FAQ~\cite{biergaizi29} also advises macOS users to
build openEMS with GCC rather than clang. The stock baseline of this report is the default clang build, and openEMS
was not built with GCC; in the standalone kernel at one thread GCC~16 runs the stock inner loop 1.16--1.39$\times$
faster than Apple clang and the rewritten one at about the same speed (the inner-loop item of
Section~\ref{sec:engine}). fdtdKit was also built on upstream's GPU work: the pre-merge \code{gpu-engine} branch at
commit 7bb294f (2026-10-12), whose GPU engine Sean Mollet added (commit bf30fa4, 2026-10-08) with follow-up work by
Thorsten Liebig, and whose gather and asynchronous dumps the fork uses (Section~\ref{sec:dumps}). The branch at
7bb294f alone runs the patch tutorial 1.13$\times$ faster than stock (both at their best), through its asynchronous
dump writer, so fdtdKit's 6.9$\times$ there (cg14, Table~\ref{tab:final}) includes about 1.13$\times$ from upstream; on
the notch tutorial, the GCPW line and patch x5 it runs at 0.98--1.03$\times$ stock (Section~\ref{sec:results}).
Upstream has since merged a reworked GPU engine, with a region-based gather, into master (c3a956b) and deleted the
branch; the merged engine was measured here on 2026-11-04 (below).

\paragraph{Concurrent work on openEMS speed.} In the weeks before and during this work, other contributors posted
speed work on openEMS, much of it, like this one, written with AI assistants. The list gives what their authors
posted, as of 2026-11-04: the numbers in it are their own, from other machines and models, and do not compare directly
with Table~\ref{tab:final}. When the list was written none of this work had been built or run here; later that day
two of the GPU engines were, on the six benchmark models (the paragraph after the list). The CPU pull requests, the
fork of Discussion \#492 and the SYCL backend were not built or run, but cg29 builds on three of them, credited in
Section~\ref{sec:cg29}: it takes the UPML part of \#244 by Noriaki Takashima (nori-proto), and it follows the
strict-order check of \#245 by Roel Bud\'e (unsettledroell404) and the lumped-RLC state slots indexed by the time step
that MrHighVoltage proposed in Discussion \#492.
\begin{itemize}
\item \textbf{A Metal engine with diamond blocking} (luikore, openEMS pull request \#228, opened
  2026-09-29~\cite{luikore228}), not merged. It advances a space-time diamond of four time steps per block; fdtdKit's
  Metal engine (Section~\ref{sec:gpu}) follows its tile schedule, without using any of its code. On a uniform grid
  with one dielectric box and PEC walls it reports 5\,752 and 4\,042\,\MC\ on an Apple M4 Pro (3.0 and 17.0 million
  cells) and 14\,734 and 11\,532\,\MC\ on an M5 Max, 1.9--6.3$\times$ the best multithreaded CPU run on the same
  machine. Its own unfused path ran at 2\,216\,\MC\ against 2\,455 for the CPU at 3.0 million cells on the M4 Pro, in
  line with Section~\ref{sec:ceiling}: on these machines the blocking, not the GPU as such, carries the gain. Its
  author calls the uniform grid Metal's best case and expects real PCB geometry to lose much of the speed-up. The
  engine is Cartesian only, a model with ADE materials (which include openEMS's conducting sheets,
  Section~\ref{sec:engine}) aborts on its default path, and its fields are checked against the SSE engine within
  relative-L2 limits, not bit for bit.
\item \textbf{Upstream's GPU engine} (Sean Mollet~\cite{mollet235}): the Metal and HIP engine that this work measured
  on the pre-merge branch (Section~\ref{sec:ceiling}) went into master as c3a956b on 2026-10-17. The merged version
  was measured here on 2026-11-04 (below).
\item \textbf{CPU engine changes with identical results.} Pull request \#231 (ismailakdag, opened
  2026-10-05~\cite{akdag231}) lets the multithreaded engine skip the barriers of extension hooks that have nothing to
  do, the idea behind this work's hook plan, and walks the SSE UPML loops along z; its outputs are bitwise identical
  on 24 setups, and it reports 1.2$\times$ (Mur) to 5$\times$ (PML-heavy) at one thread on a Ryzen~9 7900X and
  1.30--1.67$\times$ on an M5 Pro on an earlier base. A fork by MrHighVoltage (Discussion \#492,
  2026-10-26~\cite{mrhv492}) adds an AVX2 engine (x86 only), trapezoidal temporal blocking on the CPU and the GPU,
  barrier elision (552 to 189 barriers per time step), a thread tuner that scores real time steps, an asynchronous
  dump writer and a Vulkan GPU engine; it reports 1.7--4.45$\times$ multithreaded against upstream's SSE engine on
  x86 machines, with results bit-identical or differing by float32 rounding. Smaller pull requests multithread the
  Lorentz and Debye updates (\#245~\cite{pr245}) and the UPML and Debye coefficient setup (\#244~\cite{pr244}), both
  near items of Section~\ref{sec:engine}, and shorten the mesh-delta lookup (\#230~\cite{pr230}). A SYCL backend for
  Intel GPUs (\#243~\cite{sycl243}) fuses the voltage and current steps where the extensions allow.
\item \textbf{Upstream's plan} (Thorsten Liebig, Discussion \#495, 2026-10-31~\cite{thliebig495}): first a
  refactoring that separates the mesh geometry from the coefficient storage, with the speed pull requests on hold
  until it is done; then GPU kernels written once and shared by all backends. Fused voltage and current steps,
  asynchronous dump writing and temporal blocking are to go into that shared code once, temporal blocking last. Speed
  claims are to be measured with the \code{benchmark} tool of master against the existing engine on the same machine.
\end{itemize}

\paragraph{Two of these engines on the M5.} On 2026-11-04 upstream master at 5fecbb5 (2026-11-02), with its CPU engine
and \code{--engine=gpu} (the Metal backend of c3a956b), and luikore's pull request at 55ce11b, with its diamond path
and its legacy two-dispatch path (\code{OPENEMS\_METAL\_FUSED\_PIPELINE=0}), were built and run on the six benchmark
models next to stock and cg18 (\code{results/others\_2026-11-04/README.md}). The runs used this report's
\code{bench.py} (fixed step counts, the faster of two runs, port probes and $S$-parameters compared with stock), not
master's \code{benchmark} tool that Discussion \#495 asks for. Other work shared the Mac (a KiCad render loop from
another session at about 950\,\% CPU in bursts, Spotlight indexing), the other processes' GPU use was not measured,
and from about 06:00 another session was developing fdtdKit's Metal engine and may have used the GPU in the later runs,
so the numbers are rough. The SMA launch there has the 0.871\,M-cell mesh of the toolkit's current default, not the
1.071\,M cells of Table~\ref{tab:final}. luikore's diamond path refused all six models before the first step, with the
message that the Mur extension (on the GCPW line the conducting-sheet extension) ``has not migrated to the diamond
wavefront''; its legacy path, which that message reserves for diagnostic runs, ran them at 172--884\,\MC, below cg18
on every model. With every Mur face replaced by PEC (speed-only test cases, whose $S$-parameters are not those of the
real models; the GCPW line is left out, its sheets being refused too), the diamond path ran patch x5 and notch x3 at
3\,128 and 3\,734\,\MC, 2.00 and 1.86 times cg18's CPU engine, the notch tutorial at 1.47, the SMA launch at 1.05 and
the patch tutorial at 0.39 times, and at 2.1 to 6.4 times its own legacy path: on the M5, too, the blocking carries
the gain. Master's GPU engine ran the six models as they are at 0.42 to 1.06 times cg18: about level on patch x5
(1\,529 against 1\,510\,\MC) and on the SMA launch (1\,084 against 1\,019), 0.82 times on the GCPW line, and slowest
relative to cg18 on the patch tutorial (782 against 1\,856). On the PEC variant of the notch tutorial (0.270\,M cells)
it ran at 2\,374\,\MC\ against cg18's 2\,053, the faster of its two runs (1.82\,s, other processes at 43\,\% of a core
in the second before it; the other run at 1\,452\,\MC;
\code{notch\_s1\_n16000\_master-nomur.txt} in that folder), above the roughly 2\,000\,\MC\ that 74\,B per cell allow
from the M5's memory (Section~\ref{sec:background}); whether its kernels fuse the two half steps or this small model
stays in the GPU's caches was not examined. cg18 was identical to stock in every run, and so were
master's and luikore's CPU engines, except master's on the SMA launch (probes $1.5\times10^{-4}$ relative, $S_{11}$
$3.4\times10^{-3}$\,dB), a change in master since v0.37.0 that its GPU engine shows too. Apart from that, the GPU
engines differed from stock by $1.7\times10^{-7}$ to $1.3\times10^{-6}$ in the probes and by at most
$2.6\times10^{-4}$\,dB in the $S$-parameters (notch tutorial; within $3\times10^{-5}$\,dB elsewhere): master's on five
models by the same printed figures as fdtdKit's Metal engine (Section~\ref{sec:policy}), and luikore's two paths, on
the PEC variants, by the same figures as each other. None of them was compared with a CPU run with
\code{OPENEMS\_FTZ=1}. fdtdKit's Metal engine was not run on the PEC variants, and luikore's diamond path refused the
models as they are (its legacy path ran them), so on these models the two diamond engines have not been compared head
to head. They were on 2026-11-11, on luikore's own benchmark model and with master's engines beside them
(Section~\ref{sec:threebuilds}, Table~\ref{tab:threebuilds}): fdtdKit's Metal engine (cg31) ran 3\,981 and
3\,309\,\MC\ at 3.0 and 17.0 million cells, luikore's 3\,952 and 2\,829 and master's \code{gpu} engine 1\,638 and
1\,566.

Against this background the parts of fdtdKit that we have not found elsewhere are narrower than its speed-ups: the
acceptance rule applied throughout (bit identity on six models, with the closed-form validation suite and its
negative controls; Sections~\ref{sec:method} and~\ref{sec:repro}), the allocation-placement valley of the macOS
allocator, which slows stock openEMS too (Section~\ref{sec:valley}), the measured ceiling of one base-M5 memory bus
for the CPU and for the unfused GPU engine of the fork's base (Section~\ref{sec:ceiling}), with a fused and a
diamond-blocked Metal sweep
measured against it, and a diamond-blocked Metal engine that runs all six models on the GPU with every extension they
use (Mur and UPML boundaries, lumped ports, conducting sheets; Drude/Lorentz and Debye materials as well), bit-identical
to the CPU engine run with \code{OPENEMS\_FTZ=1} (Section~\ref{sec:gpu}).

\paragraph{Where fdtdKit stands.} For printed-circuit-board RF models on the CPU of Apple silicon, fdtdKit is to our
knowledge the fastest configuration of openEMS reported that keeps stock's results bit for bit: 3.2--6.9$\times$ stock
on the four models of up to about a million cells with snapshot cg14 (Table~\ref{tab:final}, the SMA launch on its
1.071\,M-cell mesh; Table~\ref{tab:current} has cg29 and Table~\ref{tab:cg31} the installed build, cg31), and,
with the modelling levers of Section~\ref{sec:modelling}, 15$\times$ end to end on the
SMA launch with the toolkit's defaults since 2026-11-04 (the 0.05\,mm mesh floor and a hold of 0.5, against stock's
full run on the 0.02\,mm mesh; cg18, Table~\ref{tab:e2e}), or 5.8$\times$ with the toolkit's defaults until
2026-11-04 (0.02\,mm and a hold of 1). On the GPU its opt-in Metal engine keeps the CPU engine's results bit for bit
when the CPU run also flushes subnormals, runs all six models on the GPU with every extension they use
(Section~\ref{sec:gpu}), and runs notch x3 and patch x5 at 3\,226 and 2\,917\,\MC\ (cg28, medians of five paired
rounds on a quiet Mac: other processes on 0.5--1.3 cores and 2--4\,\% of the GPU;
\code{results/diamond\_engine\_cg28\_tuner\_2026-11-04.txt}), where cg18's CPU engine reached 2\,046 and 1\,716 at
best on 2026-11-03 (\code{results/*\_cg18.txt}); with cg31 it runs them at 3\,702 and 3\,100
(Table~\ref{tab:cg31}). Of the other GPU engines measured on the M5, luikore's diamond path
refused all six models and master's ran them at 0.42 to 1.06 times cg18's CPU engine; on luikore's own benchmark
model cg31's Metal engine was level with luikore's at 3.0 million cells, 1.17$\times$ as fast at 17 million, and
2.1--2.4$\times$ as fast as master's \code{gpu} engine (Table~\ref{tab:threebuilds}). Whether diamond blocking on the
GPU keeps its gain on a base M5 and on PCB models is answered for fdtdKit's engine: on the M5 the fused, unblocked
Metal sweep, with its best tile, ran notch x3 and patch x5 at 0.58 and 0.65 times the blocked one (paired runs, other
processes on 2.5--4 cores) and the SMA launch (0.871\,M cells) and the GCPW line at 0.74 and 0.77 times (the faster
of two unpaired runs, other processes
on 3.3--5.6 cores; \code{results/fused\_sweep\_k1\_2026-11-04.txt}).

What is not shown: luikore's reported 4\,042--14\,734\,\MC\ on the M4 Pro and the M5 Max have no like-for-like fdtdKit
figure, because fdtdKit was not run on luikore's machines (on luikore's model it was, on this M5:
Table~\ref{tab:threebuilds}), and on the six models its diamond path ran here only on the PEC variants, next to
cg18's CPU engine and not next to fdtdKit's Metal engine; and master's GPU engine, about level with cg18's CPU engine
on patch x5 and the SMA launch, was not run on the six models next to cg29, cg31 or fdtdKit's Metal engine (the same
table has it on luikore's model). fdtdKit closes none of the
functional gaps listed above; closing some of them (Cartesian sub-gridding, conformal conductors, automatic mesh
convergence, PCB import) is the natural continuation and would need a different acceptance rule than bit identity,
since stock openEMS has no reference to be identical to.

\subsection{Outlook: clusters of Apple-silicon machines}\label{sec:cluster}
This subsection is an outlook: no cluster was built or run. Apart from the memory figures, which were measured for
this report on small boxes, and one STREAM figure of the M5 Pro (Section~\ref{sec:m5pro}), its numbers are published
hardware figures or arithmetic on them. The result of
Section~\ref{sec:ceiling} (on one machine the CPU engine is limited by one memory bus) points to the natural way to
scale: add buses. Every Mac mini or Mac Studio brings its own unified memory and its own bus, and since macOS 26.2
machines connected by Thunderbolt~5 can exchange data by remote direct memory access (RDMA)~\cite{rdmatb5}, over
links with a nominal rate of 80\,Gb/s (real throughput is lower). Published figures for current machines: a Mac mini
with M5 Pro has up to 64\,GB at 307\,GB/s~\cite{macmini}; a Mac Studio with M5 Ultra up to 512\,GB at
1.2\,TB/s~\cite{studiom5}. On a 24\,GB Mac mini with M5 Pro, STREAM's Triad measured 259\,GB/s at 15 threads, and the
five Super cores alone reached that limit (Section~\ref{sec:m5pro}).

\paragraph{Parameter studies.} Two simulations side by side on one machine gain nothing, because they share the bus
(Section~\ref{sec:rejected}). On separate machines they do not, so a sweep or an optimisation run spread over $n$
machines should scale close to $n$ with no change to the engine, only a job dispatcher. This is the cheapest use of a
cluster and the one most RF design work needs.

\paragraph{One large model.} openEMS had an MPI engine for domain decomposition up to v0.0.36, but it had not compiled
for years (it used the C++ MPI bindings that MPI-3 removed), and upstream removed it before v0.37.0 (commit 5261463,
2026-09-16). A model split across machines therefore needs a new domain-decomposed engine, which could carry the
fork's engine changes from the start. One fork has tried an x-slab split of one model across two GPUs with
host-staged halo exchange (td-krzysiek, described in Discussion \#495~\cite{thliebig495}); upstream's maintainer
replied there that multi-device engines are not a priority and that optimising the mesh usually pays more.
The arithmetic is favourable. During the time loop a plain cell holds about 28 bytes (six float field components and a
4-byte operator index), but the operator setup peaks near 100 bytes per cell, about 124 with a field-dump or nf2ff
box, and on a Mac the two share the same memory (measured for this report with snapshot cg18 at 2 threads on boxes of
up to $64^3$ cells; \code{results/memory\_per\_cell\_2026-11-03.txt}). cg29 allocates the operator index once per
z-vector of four cells instead of once per cell, about 3 bytes per cell less (patch 0017; not measured on these
boxes). Memory alone would therefore allow about four to five billion cells on one 512\,GB Mac Studio, before the
operating system's share; but the arrays use 32-bit sizes, and the code refuses meshes above about 4.29 billion nodes,
or 1.43 billion with field dumps, an nf2ff box, conducting sheets or the basic engine or \code{--engine=gpu};
\code{--engine=metal-diamond} keeps a model on the GPU only below $2^{32}$ floats per field array (about 1.43 billion
nodes) and up to 65\,535 lines in x or y, and runs the one-thread CPU loops above that (Section~\ref{sec:gpu}). A setup
that does not hold the whole mesh at once, and 64-bit sizes, would move the limit towards the run state. At equal
bytes per cell, 512\,GB is 6.4 times the 80\,GB of an H100. With the domain split into slabs, each half step exchanges
one plane of tangential fields per neighbour: for a $1000^3$-cell model on four machines that is about 8\,MB per
interface and half step, 16\,MB per step, about 1.6\,ms at the nominal link rate (more at the real throughput),
against roughly ten milliseconds of update per step if each machine ran at the bandwidth roofline. With the exchange
overlapped with the interior update, four M5 Ultra machines would offer an aggregate bus of about 4.8\,TB/s, above the
3.35\,TB/s of one H100~\cite{h100}. Two caveats: whether the CPU cores of an Ultra-class chip can draw on the whole
bus, and how far the Metal engine gets there, have to be measured (on the M5 a fused, unblocked Metal sweep stayed
near the fused CPU sweep, and only the diamond-blocked one ran the large models clearly faster,
Section~\ref{sec:gpu}); and the bit-identity rule carries over naturally, since a domain decomposition need not change
any cell's arithmetic, so a distributed run can be checked against a single-machine run to the bit.
